
\section{Related Work}
\label{sec:related}

\noindent
Related work includes findings in the fields of visualization of tree-structured and hierarchical data, non-photorealistic rendering (NPR), human perception, map design, and cartographic information display. 

Since the original presentation of Johnsons and Shneiderman's 2D treemaps \cite{Johnson1991} 
various treemap layouts have been investigated, 
improving layout readability, stability, and the graphical elements' aspect ratio \cite{Bruls1999,Bederson2002}.
2.5D treemaps utilize the third dimension for additional information display based on the Cityscape metaphor  \cite{Keskin1997} by transforming rectangular elements into stacked 3D blocks.  
Our GPU-based implementation follows the  real-time rendering approach for 3D treemaps described in \cite	{Trapp2013} that can efficiently process large amounts of tree data. 

Sketchy rendering is often based on perturbing outlines (\eg silhouettes) and modifying textures (\eg hatches, strokes) in a sketchy way \cite{Strothotte2002}.
For example, if we map uncertainty values to edges, drawn in  image space \cite{Nienhaus2004}, we get sketchy outlines causing a kind of \enquote{shower door} effect. 
Alternatively, overlapping, edge-aligned and textured billboards can be placed in object space, 
enabling image-space independent sketchiness \cite{Doellner2003}.
In this case, efficient and accurate line visibility can be achieved by GPU-based hidden line removal \cite{Cole2009}.
In this work, we follow the object-based approach for sketchy rendering. 
To draw textures in a sketchy style, we can hatch strokes over surfaces, \eg via precomputed \cite{Praun2001,Webb2002} or recursive procedural Tonal Art Maps \cite{Szecsi2014}. 
For hatching, a GPU-based, dynamic principal direction propagation emphasizes surface orientation \cite{Kim2008} but does not significantly contribute to sketchiness of blocks in 2.5D treemaps due to their simple, regular shape. 
Most computer generated NPR styles come close to hand-drawn ones and seem to be preferred by users \cite{Cole2008, Isenberg2006}.

Sketchiness in general allows us to influence how people read and understand depicted information \cite{Isenberg2013}.
Wood \etal \cite{Wood2012} investigate sketchy rendering for information visualization and proposes a corresponding framework; we extend this work for the specific application of sketchiness in 2.5D treemaps. 
In particular, sketchiness can intentionally facilitate communication \cite{Walny2011} or 
reduce the expected amount of cognitive processing due to its supposed simplicity \cite{Wood2012}.
It conveys visual imprecision, decreases people's confidence in the underlying data quality and encourages constructive feedback, discussion, and participation \cite{Zanola2009,Luboschik2010} as well as supports interactive exploration \cite{Greenberg2011, Browne2011}.

In cartographic information display \cite{Bertin1983} 
and information visualization,
visual variables  provide a theoretical background for designing and understanding visual mappings \cite{Skupin2000}.
Each visual variable (\eg size, shape, color, texture) can be used to map one of the underlying multi-dimensional data 
but its choice depends on the variable's characteristics including
selective, associative, quantitative, length, and order properties \cite{Carpendale2003}, which significantly have impact on perception and understanding of the presented information. 

Sketchiness can be viewed as visual variable 
that alters the way other visual variables are rendered, 
\ie it perturbes their visual mappings.
In particular, sketchiness appears to \enquote{have the capacity to carry information in its own right} \cite{Wood2012} 
and allows us to map any ratio-scale data, while being most effective for ordinal data. 
For example, to depict uncertainty, Boukhelifa \etal \cite{Boukhelifa2012} propose to use only 3 to 4 discrete line styles.

